\clearpage
\pdfbookmark[2]{6. Gauss an Olbers, 13. Januar 1815}{parabolisch-006}
\begin{center}
[6.]\par
\medskip
\textsc{Gauss} an \textsc{Olbers}, Göttingen, 13. Januar 1815.
\end{center}
\noindent\hspace*{\parindent}.\ .\ .\ .\ .\ .\par
Ich habe mich diese Woche hindurch noch mit der Theorie der Cometenbahnen beschäftigt, und bin noch auf eine Vereinfachung gekommen, nach welcher mir nun nichts mehr zu wünschen übrig zu bleiben scheint. Ich finde nemlich für die Zwischenzeit den Ausdruck:
\[
t''-t=3mk\BareRoot{(r+r'')}\cdot\BareRoot{\left(1-\frac{27m^4k^6}{4(t''-t)^4}\right)}\mathpunct{\raisebox{.45ex}{,}}
\]
der in allen in der Praxis vorkommenden Fällen genau genug ist. Nach völliger Strenge sollte die letzte Quadratwurzel eine unendliche Reihe sein, die, wenn ich der Kürze halber \(\frac{27m^4k^6}{4(t''-t)^4}=x\) setze, nach meiner Entwickelung wird \(=\BareRoot{(1-x-2x^3-8x^4-\ldots)}\) und wo also glücklicherweise das zweite Glied fehlt. Selbst, wenn \(k=\tfrac12(r+r'')\), würde \(x\) nur etwa \(\frac1{47}\), also der Fehler, \SourceLeaf{342}wenn man sich auf \(\BareRoot{(1-x)}\) einschränkt, nur etwa \(\frac1{100000}\) des Ganzen\HistoricalNoteEight{In einem Exemplar der deutschen Übersetzung der Abhandlung »Observationes Cometae secundi A. 1813 etc.« findet sich die handschriftliche Bemerkung:] NB. Die Reihe \(1-z-2z^3-8z^4-45z^5-272z^6-\ldots\) convergirt, so lange \(z<\frac{27}{256}\); für diese Grenze, der \(k=r+r''\) entspricht, wird ihr Werth \(=\frac89\).}. Ich wende nun obige Formel auf folgende Art an. Ich setze
\[
u=A\operatorname{tang}Q,
\]
wodurch
\[
k=\frac A{\cos Q}\quad\text{und}\quad\rho=\frac{g\cos(Q-\varphi)}{h\cos Q}\mathpunct{\cdot}
\]

Man hat also
\[
\begin{aligned}
(t''-t)^2\cos Q^2={}&9mmAA\left(1-\frac{27m^4A^6}{4(t''-t)^4\cos Q^6}\right)\\[4pt]
&\quad\Biggl\{\BareRoot{\left(BB+\left(\frac{A\operatorname{tang}Q+c}{b}\right)^2\right)}\\[4pt]
&\qquad+\BareRoot{\left(B''B''+\frac{A\operatorname{tang}Q+c''}{b''}\right)^2}\Biggr\}
\end{aligned}
\]
oder, indem ich folgende Bezeichnungen einführe
\[
\begin{gathered}
\frac{9mmAAB}{(t''-t)^2}=\delta,\qquad\frac{9mmAAB''}{(t''-t)^2}=\delta'',\\[5pt]
\frac A{bB}=\epsilon,\qquad\frac A{b''B''}=\epsilon'',\\[5pt]
\frac c{bB}=\eta,\qquad\frac{c''}{b''B''}=\eta'',\qquad\frac{27m^4A^6}{4(t''-t)^4}=\omega,
\end{gathered}
\]
wird:
\[
\begin{aligned}
\frac{\cos Q^2}{1-\dfrac\omega{\cos Q^6}}={}&\delta\BareRoot{\bigl(1+(\epsilon\operatorname{tang}Q+\eta)^2\bigr)}\\[4pt]
&+\delta''\BareRoot{\bigl(1+(\epsilon''\operatorname{tang}Q+\eta'')^2\bigr)}.
\end{aligned}
\]

Etwas einfacheres ist wohl nicht zu wünschen. Man könnte noch zwei Hülfswinkel \(\vartheta\), \(\vartheta''\) einführen, indem man
\[
\begin{aligned}
\delta\epsilon&=\lambda\sin\vartheta, &\delta''\epsilon''&=\lambda''\sin\vartheta''\\
\delta\eta&=\lambda\cos\vartheta, &\delta''\eta''&=\lambda''\cos\vartheta''
\end{aligned}
\]
setzte, wodurch die Gleichung würde
\[
\begin{aligned}
\frac{\cos Q^3}{1-\dfrac\omega{\cos Q^6}}={}&\BareRoot{\bigl(\delta\delta\cos Q^2+\lambda\lambda\cos(Q-\vartheta)^2\bigr)}\\[4pt]
&+\BareRoot{\bigl(\delta''\delta''\cos Q^2+\lambda''\lambda''\cos(Q-\vartheta'')^2\bigr)}.
\end{aligned}
\]

Man könnte selbst noch weiter gehen und durch bekannte Verwandlungen der Gleichung folgende Form geben:
\[
\begin{aligned}
\frac{\cos Q^3}{1-\dfrac\omega{\cos Q^6}}={}&\mu\BareRoot{\bigl(1+\nu\nu\cos(Q+\pi)^2\bigr)}\\[4pt]
&+\mu''\BareRoot{\bigl(1+\nu''\nu''\cos(Q+\pi'')^2\bigr)}.
\end{aligned}
\]

\SourceLeaf{343}Allein der Gewinn\footnote{Er würde nemlich \textls[150]{bloss} darin bestehen, dass man sogleich das Maximum und Minimum des Werthes der beiden Wurzelgrössen beurtheilen könnte, wodurch man, anfangs \(\omega\) vernachlässigend, \textls[150]{Grenzen} für \(\cos Q^3\) und folglich für \(Q\) erhielte. Bei der Berechnung jedes einzelnen Versuchs werden doch immer 7 Aufschlagungen erfordert, man wähle welche Form man wolle.} würde zu unerheblich sein, um die Mühe dieser Verwandlungen zu belohnen. Immer ist die Gleichung vollkommen strenge, sobald man statt des Nenners im ersten Theile die Reihe setzt
\[
1-\frac\omega{\cos Q^6}-\frac{2\omega^3}{\cos Q^{18}}-\frac{8\omega^4}{\cos Q^{24}}-\text{etc.}
\]

Das Problem ist also auf Eine Gleichung mit sieben bekannten Grössen gebracht, und ich glaube nicht, dass es möglich ist, es auf eine geringere Anzahl zu reduciren.

Die allgemeine Aufgabe, nemlich, wenn der mittelste Ort unvollständig bekannt ist, und man bloss irgend einen grössten Kreis kennt, worin er liegt — von welcher das Gegenwärtige eigentlich nur der einzelne Fall ist, wo dieser grösste Kreis durch den mittlern \(\odot\)ort geht — zu welcher man seine Zuflucht nehmen muss, so oft dieser \(\odot\)ort nahe bei der Richtung der geoc[entrischen] Bewegung liegt, diese allgemeine Aufgabe habe ich mir nun auch auf eine analoge Art durchgeführt und ein Musterbeispiel sorgfältig berechnet\HistoricalNoteEight{Siehe die folgende Notiz.]}. Der Natur der Sache nach ist die Arbeit hier viel verwickelter, ich komme am Ende auf zwei Gleichungen mit 11\HistoricalNoteDoubleUnbracketed{Jeder Versuch 11 Aufschlagungen erfordernd.} bekannten Grössen, und zweifle auch, ob es möglich ist, diese Anzahl kleiner zu machen. Wenn nach allen Abkürzungen jetzt die Berechnung einer Cometenbahn im ersten Fall im Minimum 1 Stunde erfordert (so ist es etwa bei mir), so möchte der andere Fall leicht 2—3 Stunden erfordern.
\par\medskip\noindent\makebox[\linewidth]{\rule{.18\linewidth}{.35pt}}\par
